\documentclass[a4paper,11pt]{ctexbook}
\usepackage{finance-style}

\title{\textcolor{themeblue}{\Huge\textbf{金融创新与系统性风险}}}
\subtitle{\textcolor{gray}{第三卷\quad 共十课}}
\author{}
\date{}

\begin{document}
\maketitle
\thispagestyle{empty}
\tableofcontents
\newpage

% ============================================================
\chapter{影子银行——监管之外的大玩家}
% ============================================================

\begin{scenario}
\textbf{包商银行——当"同业存单"不再刚性兑付}

2019年5月，中国人民银行和银保监会宣布包商银行被接管——这是中国20年来第一家被接管的商业银行。包商银行的资产端表外理财规模超过5000亿，是其存款规模的2倍以上。它通过发行同业存单从银行间市场大量融资，然后通过"同业理财"将资金投向非标资产——形成了"同业存单→同业理财→委外投资"的影子信贷链条。当包商银行的大股东"明天系"坍塌时，这条链条断裂了。包商银行案例揭示了中国影子银行体系的典型风险结构。
\end{scenario}

\begin{qualitative}
\textbf{什么是影子银行？}

金融稳定理事会（FSB）的定义：信贷中介活动中那些在常规银行体系之外的、存在监管套利的部分——核心特征：期限转换（借短投长）、流动性转换、高杠杆、监管缺席。中国的影子银行在2017年巅峰期接近100万亿人民币——大到"不能忽视"。

\textbf{主要形式：}
\begin{itemize}
  \item 银行理财（巅峰期约30万亿）——表外、期限错配、刚兑预期
  \item 信托贷款——房地产和平台公司的主要融资渠道
  \item 委托贷款——企业间直接借贷，银行充当通道
  \item 未贴现银行承兑汇票——表外信贷的隐蔽形式
  \item 同业存单+同业理财——中小银行从大行"批发"资金
\end{itemize}
\end{qualitative}

\begin{quantitative}
\textbf{资管新规（2018）——影子银行的大规模出清}

2017年底影子银行规模约100万亿。资管新规的核心条款：
\begin{itemize}
  \item 打破刚兑——理财不得承诺保本保息
  \item 净值化管理——理财必须按市值计价
  \item 消除多层嵌套——禁止资管产品投资资管产品
  \item 限制非标——理财投资非标资产不得超过净资产的35\%
\end{itemize}

到2024年，影子银行规模压缩至约60万亿——减少了40\%。

\textbf{资金池模式的数学崩溃：}资金池模式下，银行用一个"大池子"吸收短期理财资金、投向长期非标资产。只要池子持续有净流入（新理财资金 > 赎回），银行可以用新钱付旧钱——但一旦净流入转为净流出（赎回>新发），池子就面临流动性危机。崩盘速度取决于：非标资产的到期期限（3-5年）vs 理财产品的期限（3-6个月）的错配程度。
\end{quantitative}

\begin{lessonsummary}
\begin{enumerate}
  \item 影子银行是银行为规避资本充足率和存贷比监管而进行的"监管套利"
  \item 中国影子银行的核心风险——期限错配+刚兑预期+多层嵌套
  \item 2018年资管新规显著降低了影子银行的规模，但非标融资的真空尚未被有效填补
  \item 影子银行不是"坏"的——它提高了信贷供给效率——问题在于其风险没有被充分定价和监管
\end{enumerate}
\end{lessonsummary}

\begin{discuss}
\begin{enumerate}
  \item 影子银行收缩的40万亿去哪了？——一部分回了表内（成为贷款），一部分变成了直接融资（债券），一部分消失了（依赖非标融资的房地产和城投平台失去了融资来源）。这直接解释了为什么2018-2020年很多民营企业出现了"融资难、融资贵"的问题。
  \item 美国版的影子银行——货币市场基金（MMF）、私募信贷基金、回购市场（Repo）——在2020年3月的COVID市场崩溃中也出现了类似的资金池崩塌。美联储被迫推出了货币市场共同基金流动性便利（MMLF）——央行第一次直接为影子银行体系提供流动性。为什么央行不救"影子银行"？——因为如果不救，影子银行的崩溃会通过信贷收缩传递到实体经济。央行面临"两害取其轻"的抉择。
  \item 更好的监管方案不是"禁止"影子银行——而是让它"见光"——要求注册、披露、计提资本。中国的做法是"监管全覆盖"——将所有资产管理产品纳入统一监管框架。这个方向对吗？——对，但执行中的挑战是人力不足（监管人员的数量vs被监管机构的数量完全不匹配）。
\end{enumerate}
\end{discuss}

\begin{exercises}
\begin{enumerate}
  \item 一个银行理财资金池——吸收3个月期限的理财产品（利率3\%）、投资3年期的非标资产（利率6\%）。如果每月新发理财10亿、赎回8亿（净流入2亿），在净流入持续的情况下，池子可以维持多久？如果市场恐慌导致净流出变为-2亿（赎回10亿、新发8亿）——资金池的流动性缺口在第几个月变得不可承受？
  \item 假设你持有1000万的某银行理财（底层资产是非标城投债）。银行告诉你"该理财净值从1.0降到了0.85"——你损失了15\%。但如果银行执行的是"成本法"估值（不是市值法）——这些损失可能在账面上根本看不到。用成本法估值的"平滑效应"测算：当底层非标资产实际损失了20\%时，成本法估值的理财净值可能只显示损失了多少？
  \item 中国影子银行规模的官方统计口径（社融中的"表外三项"）和非官方口径（包含"银行理财+银信合作+券商资管"）的差距是多少？查一下最新的数据并分析原因。
\end{enumerate}
\end{exercises}

\xuehaiinline{
\textbf{中国的"同业存单—同业理财—委外"链条解析：}（1）A银行（通常是股份制或大型银行）在银行间市场发行同业存单（CD），吸收B银行（中小银行）的同业资金——B银行的存款付息率1.5\%，同业存单收益率3\%+，B银行有动力把存款换成同业存单；（2）B银行再将从CD募集的资金买入C银行发行的同业理财产品——收益率再加100-200bp；（3）C银行将同业理财的资金委托给D券商/基金子公司做"委外投资"——投向加了杠杆的债券或非标。这个链条在2016年达到顶峰——B银行用拆入的同业资金买了C银行的同业理财，C银行用这笔钱做了委外投资——本质上是银行体系内部的钱在"空转"，没有流入实体经济。2017年央行把同业存单纳入MPA考核（宏观审慎评估体系）后，这一链条开始收缩。
}

\newpage
% ============================================================
\chapter{2008年全球金融危机——深度解剖}
% ============================================================

\begin{scenario}
\textbf{雷曼兄弟倒下的48小时——全球金融海啸的起点}

2008年9月13-15日，周末。美国财政部、美联储和华尔街最大的银行CEO们在纽约联储大楼紧急开会——讨论是否救助雷曼兄弟。周末结束时，没有买家出现（美国银行和巴克莱都放弃收购），政府也拒绝救助。9月15日凌晨，雷曼——158年历史的全球第四大投资银行——申请破产保护。道琼斯指数当天暴跌超500点。全球信贷市场在48小时内冻结。但雷曼的资产并非"一文不值"——它的净资产在破产申请日仍然为正（约260亿美元）。问题在于：流动性枯竭了——没有人在恐慌的市场中买任何有风险的东西。
\end{scenario}

\begin{qualitative}
\textbf{危机的七步链}

\textbf{第一步——低利率环境：}2000年互联网泡沫破裂后，美联储将联邦基金利率从6.5\%降至1.0\%（2003-2004年）。极低利率催生了房地产泡沫。

\textbf{第二步——房贷泡沫：}银行大量发放房贷——甚至向没有稳定收入的人发放（"次贷"）。2006年次贷占新发放房贷的约20\%。

\textbf{第三步——证券化狂热：}这些次贷被证券化为RMBS。RMBS再被分层打包成CDO。CDO的CDS再被打包成合成CDO。基础资产和终端产品之间的距离越来越远——没有人知道底层是什么。

\textbf{第四步——评级腐败：}标准普尔和穆迪给大量MBS和CDO打出了AAA评级——和国债同级别。2007年标准普尔收入的44\%来自结构化金融产品的评级——存在系统性的利益冲突。

\textbf{第五步——杠杆堆积：}华尔街的五大投行的平均杠杆率从2003年的约20倍升至2007年的约30倍（Bear Stearns达到33倍）。30倍杠杆意味着资产价格仅需下跌3.3\%就会清零股权。

\textbf{第六步——房价转跌：}2006年美国房价见顶。2007年开始下跌。失去了"房价永远涨"的信仰后，次贷违约率飙升。

\textbf{第七步——多米诺坍塌：}贝尔斯登被收购→房地美和房利美被接管→雷曼破产→AIG被救助→美林被收购→高盛和摩根士丹利转为银行控股公司→全球信贷冻结。
\end{qualitative}

\begin{quantitative}
\textbf{通过资产负债表传导的传染模型}

假设三个银行A、B、C之间有同业存款关系：A存了100亿在B，B存了100亿在C。如果A倒闭，B可能无法拿回在A的存款——B的资产负债表受损——如果受损后的B无法偿还C的负债——C也受损。这就是\textbf{资产负债表传染}。

\textbf{Eisenberg-Noe模型（2001）：}给定银行间风险暴露矩阵$X$（$X_{ij}$表示银行$i$欠银行$j$的金额），以及每个银行的外部资产和外部负债，可以求解在清算时每家银行能回收的比例。模型揭示了一个令人不安的性质：当网络足够互联时，一个中小银行的违约可以通过"多米诺效应"击倒数倍于自身规模的大型银行。
\end{quantitative}

\begin{lessonsummary}
\begin{enumerate}
  \item 2008年金融危机不是单一原因导致的——它是低利率、杠杆堆积、监管失败、评级腐败、风险模型缺陷等多重因素的叠加结果
  \item 银行间市场的互联性创造了"资产负债表传染"渠道——一个小银行的违约可以蔓延到整个系统
  \item "大而不能倒"（Too Big to Fail）——政府面临两难：救助滋生道德风险，不救助引发系统性崩溃
  \item Dodd-Frank法案（2010）试图解决"大而不能倒"——但2023年的SVB危机说明问题远未被解决
\end{enumerate}
\end{lessonsummary}

\begin{discuss}
\begin{enumerate}
  \item Greenspan在危机后承认："我发现了我的世界观中的一个缺陷——我以为市场的自利动机足以保护股东和存款人——但我错了。"这句话意味着什么？——连做了20年联储主席的人对自己的信仰产生了怀疑，说明了什么？
  \item "系统性风险"和"单个机构风险"的本质区别是什么？——单个机构的风控是"我怎么不死"；系统性风险的分析是"你怎么死会连累到我"——两者需要完全不同的分析框架。
  \item 中国是否面临类似2008年美国式的系统性风险？差别在哪里？——中国的银行体系以国有大行为主体、资本市场规模相对较小、资本账户不完全开放——这些结构性差异使"中国版次贷危机"发生的概率较低。但"不一样"不等于"没有风险"——中国的风险集中在房地产和地方债务，而不是风险证券化。
\end{enumerate}
\end{discuss}

\begin{exercises}
\begin{enumerate}
  \item 构建一个3银行的简化模型：A银行有100亿外部资产和80亿外部负债；B银行有80亿外部资产和60亿外部负债；C银行有120亿外部资产和100亿外部负债。银行间风险暴露：A存了20亿在B，B存了15亿在C，C存了10亿在A。如果A的外部资产价值下跌了30\%——计算A的净值变化，以及通过同业暴露传导后B和C的净值变化。
  \item 分析"顺周期性"（Procyclicality）——为什么银行的杠杆在经济好的时候自动上升、在经济差的时候被迫下降？这种顺周期性如何放大了经济波动？（提示：VaR模型在好年景计算出来的风险小——银行可以加更大杠杆——但这正好在泡沫时期积累了最大的隐患。）
  \item 下载2003-2012年美国大型银行（如JPMorgan, Citigroup, Bank of America）的季度杠杆率数据，画出时间序列。标注2008年9月的位置。在雷曼破产前，这些银行的杠杆率在什么水平？
\end{enumerate}
\end{exercises}

\xuehaiinline{
\textbf{沃尔克规则（Volcker Rule）与"自营交易"的争议：}Dodd-Frank法案的沃尔克规则禁止银行用自有资金进行自营交易（Proprietary Trading）——即用银行的资本金赌市场的方向。支持者认为：自营交易是2008年危机中银行过度冒险的根源。反对者认为：自营交易在大型银行的总风险敞口中占比不大（通常不到10\%），真正危险的是银行对结构化产品的隐性敞口和杠杆。2019年美联储放松了沃尔克规则——但在2023年SVB危机后，重新收紧的呼声再起。
}

\newpage
% ============================================================
\chapter{LTCM与模型风险——当诺贝尔奖也救不了你}
% ============================================================

\begin{scenario}
\textbf{1998年——"10个标准差事件"摧毁了华尔街最聪明的人}

LTCM在1998年1月还有50亿美元资本金。到9月仅剩6亿——几周内归零。它的合伙人包括两位诺贝尔奖得主（Scholes, Merton）和美联储前副主席（David Mullins）。它的交易员是华尔街最优秀的债券套利专家。它的风控模型是当时最先进的。但市场的一个"10个标准差事件"——按照俄罗斯违约的概率在LTCM的历史数据中几乎为零——摧毁了一切。索罗斯的评论："这次看到的是一个数学灾难——即使知道概率，模型也会在人们最需要它的时候失效。"
\end{scenario}

\begin{qualitative}
\textbf{LTCM的收敛交易策略}

LTCM的核心策略——收敛交易（Convergence Trading）：识别两种高度相关的金融资产，当它们的价差（Spread）暂时偏离历史水平时，做多被低估的、做空被高估的——赌价差回归。

\textbf{典型交易：}做多意大利国债（收益率高）、做空德国国债（收益率低）——赌意大利和德国的利差缩小（因为两国最终都会加入欧元）。这个交易在1997-1998年的大部分时间里都赚钱。

\textbf{问题：}LTCM在数十个类似的"收敛"交易上同时建仓，且每个交易的规模都极大——加上平均约30倍的杠杆（最高达到100倍）。
\end{qualitative}

\begin{quantitative}
\textbf{LTCM的杠杆毁灭计算}

假设LTCM的一个收敛交易——做多意大利债券（久期5年）、做空等面值的德国债券（久期5年），利差200bp。LTCM用了50倍杠杆（每1元资本金做50元的交易）。

如果利差仅扩大8bp（0.08\%）——两个债券价格各变动约5年×0.08\%=0.4\%——双向损失约0.8\%——乘以50倍杠杆就是40\%的资本金损失。

事实上，1998年8-9月，俄罗斯违约后，意大利-德国的利差从约200bp扩大到了约350bp——扩大了150bp。如果用上面的简化模型估算：损失≈5年×1.5\%×2倍（双向）×50倍杠杆=750\%——即资本金清零了7.5次。

\textbf{相关系数的"体制转换"：}俄罗斯违约前，LTCM的模型中所有新兴市场经济体的资产之间的相关系数设定为0.3-0.5。违约后，这些相关系数在几天内飙升到0.8-0.95——市场一致避险，"分散化投资"同时失效。
\end{quantitative}

\begin{lessonsummary}
\begin{enumerate}
  \item LTCM的核心错误不是策略错误——而是杠杆的过度使用和模型的过度信任
  \item 50倍杠杆意味着2\%的反向价格变动就会清零资本金——但市场在极端情况下可以出现10-20\%的反向变动
  \item 模型风险不是"模型本身错了"——而是"在关键的时候模型假设不成立了"
  \item "尾部风险"不能被历史数据估计——过去100年没有发生过的事，明天可能发生
\end{enumerate}
\end{lessonsummary}

\begin{discuss}
\begin{enumerate}
  \item LTCM和2008年危机之间的共性是什么？——都是"杠杆+相关性突变+流动性枯竭"的三重打击。不同的只是产品（LTCM做债券利差，2008年做CDO）——背后的数学结构完全相同。这是否意味着金融风险的根本模式是相同的，变的只是载体？
  \item "黑天鹅"（Nassim Taleb的概念）和"灰犀牛"（Michele Wucker的概念）——LTCM的崩盘是黑天鹅（不可预见的俄罗斯违约）还是灰犀牛（高杠杆的脆弱性一直明知但没人处理）？
  \item 现代金融体系是否比LTCM时代更"抗风险"了？——是的：巴塞尔协议的资本金要求提高了、中央清算机制建立了、压力测试变得常规了。但也不是——影子银行的规模更大、高速交易技术的"闪崩"风险更大了、加密货币和DeFi在监管真空中积累了巨大风险。金融风险的性质变了，但在总量上看并没有减少。
\end{enumerate}
\end{discuss}

\begin{exercises}
\begin{enumerate}
  \item 给定一个收敛交易——做多A国债（久期7年，收益率6\%）、做空B国债（久期5年，收益率4\%），名义本金1亿元。初始利差200bp。假设利差扩大到250bp、300bp、400bp——分别计算单向损失（注意：久期越长对利率越敏感，所以A的损失比利差扩大对久期的影响更复杂——这里简化处理为净久期差引起的损失）。
  \item 假设LTCM有100亿资本金，平均杠杆率30倍，组合中所有交易的隐含波动率假设为5\%年化。如果实际波动率飙升到30\%——在VaR$_{99\%}$（日度）下，一个交易日的潜在损失从多少变成了多少？
  \item 在真实历史中，俄罗斯违约后的卢布贬值了约70\%。假设你的组合中有5\%的俄罗斯资产——卢布贬值70\%加资产价格本身下跌50\%——对你的组合的总损失贡献是多少？为什么即使很小的比重也能造成重大影响？（提示：杠杆的放大效应）
\end{enumerate}
\end{exercises}

\xuehaiinline{
\textbf{Nassim Taleb的"反脆弱"思想：}Taleb在其著作《反脆弱》中提出——"黑天鹅"（极端不可预见的冲击）是金融和生活的固有特征，而不只是"异常"。因此，正确的策略不是试图预测黑天鹅（做不到）——而是构建一个在"好"的时候牺牲一些收益、在"坏"的时候能够从冲击中受益或至少不受致命损失的系统。这被称为"杠铃策略"（Barbell Strategy）——绝大部分资产放在极度安全的投资中（如国债），极小部分放在极高风险的投机上（如期权）。中间地带——LTCM的"中等风险、高杠杆"的收敛交易——是最脆弱的：可预测的日常波动是小的，但一次黑天鹅就可以清零全部。
}

\newpage
% ============================================================
\chapter{做空——使市场完整的必要之恶}
% ============================================================

\begin{scenario}
\textbf{Michael Burry——《大空头》的主角如何赌赢次贷危机}

2005年，对冲基金经理Michael Burry仔细研究了MBS（住房抵押贷款支持证券）的底层数据——发现大量房贷的借款者根本没有还款能力。但市场上几乎所有人都认为"房价永远不会跌"。Burry决定做空MBS——通过买入针对次级房贷的CDS（信用违约互换）。他的策略在2005-2007年持续亏损（每年支付保费），投资人纷纷撤资。到2008年——他的基金赚了约7亿美元。Burry为什么能赢？因为他看到了别人没有看到的东西——并且有了做空的工具。
\end{scenario}

\begin{qualitative}
\textbf{做空的本质}

做空（Short Selling）是通过借入并卖出资产、或者通过衍生品（CDS、看跌期权）来从价格下跌中获利。做空在舆论上经常被批评——"做空是在赌别人亏钱"——但它在市场中的经济功能是必不可少的：

\begin{enumerate}
  \item \textbf{价格发现：}做空者揭露估值泡沫和高估的公司——没有做空者，市场价格只会反映乐观者的看法
  \item \textbf{市场流动性：}做空者增加了市场的买卖双方平衡——使交易更容易执行
  \item \textbf{风险管理：}做空为持有现货的投资者提供了对冲工具——可以锁定利润或限制损失
\end{enumerate}

\textbf{做空的两种形式：}
\begin{itemize}
  \item \textbf{直接做空：}借股票 → 卖出 → 获得现金 → 未来买回归还。风险和收益不对称——最大收益为100\%（股价归零），最大损失是无限的（股价可以涨到无穷）
  \item \textbf{通过衍生品做空：}买入看跌期权（有限损失）或买入CDS（做空信用风险）
\end{itemize}
\end{qualitative}

\begin{quantitative}
\textbf{做空一只股票的损益计算}

借入1000股并以50元卖出→得到50000元。一个月后：

\begin{itemize}
  \item 股价跌到30元：以30000元买回1000股归还→盈利20000元（+40\%）
  \item 股价跌到10元：以10000元买回→盈利40000元（+80\%）
  \item 股价涨到70元：以70000元买回→亏损20000元（-40\%）
  \item 股价涨到100元：以100000元买回→亏损50000元（-100\%）
  \item 股价涨到200元：以200000元买回→亏损150000元（-300\%）
\end{itemize}

这显示了做空的"不对称风险"——潜在无限损失。因此做空者必须设置严格的止损。

\textbf{轧空（Short Squeeze）的数学：}当被大量做空的股票价格上涨时，做空者被迫买入平仓——这些买入进一步推高了股价。如果空头比率（Short Interest / 日均交易量）非常高，轧空的威力可能极其剧烈。GameStop在2021年1月的空头比率一度超过140\%（同一股票被多次借出做空）。
\end{quantitative}

\begin{lessonsummary}
\begin{enumerate}
  \item 做空是市场完整运行的必要组成部分——它提供了价格发现、流动性和对冲工具
  \item 做空的风险收益高度不对称——最大收益有限（100\%）、最大损失无限
  \item 轧空（Short Squeeze）是被做空股票价格快速上涨可能引发的恶性循环——做空者被迫买入→进一步推高股价
  \item 没有做空工具的市场（如中国A股在2010年之前）是不完整的——坏消息只能通过卖出持有股票来释放，效率远低于可以通过做空直接传递
\end{enumerate}
\end{lessonsummary}

\begin{discuss}
\begin{enumerate}
  \item 中国证监会2023年出台了对融券做空的限制——对融券保证金比例从50\%提高到80\%，限制战略投资者出借股票。支持的理由：做空加剧了市场下跌；反对的理由：做空是市场效率的必要工具。哪种观点更有道理？——取决于你问谁：散户投资者大多支持限制做空（认为做空者是"坏人"）；机构投资者大多反对限制（认为做空是风险管理和定价必要的工具）。
  \item 2015年A股股灾期间，中国监管机构一度禁止做空（甚至禁止大股东减持）。效果如何？——市场仍然下跌，因为禁止做空没有解决导致下跌的根源（杠杆踩踏），反而消除了市场恢复平衡的一种机制。这是"打信号弹的人"还是"引发恐慌的人"的经典争论。
  \item 如果CDS做空MBS的Burry被禁止做空——次贷危机会更晚被发现、泡沫会更大、最终的崩盘会更惨。做空者的社会价值——"唱空"是令人不快的，但可能是防止泡沫失控的唯一制衡力量。
\end{enumerate}
\end{discuss}

\begin{exercises}
\begin{enumerate}
  \item 一只股票当前价格80元。做空者以80元借入1000股卖出，融券成本年化8\%。一个月后股价涨到90元——做空者的亏损额和亏损百分比是多少？如果两个月后股价涨到120元呢？三个月后涨到200元呢？
  \item 在GameStop轧空案例中，假设做空者的平均成本20元、最高空头仓位1亿股（来源：SEC报告）。股价最高涨到483元——做空者的累计浮亏峰值为多少？（答案：约460亿美元）
  \item 做空者和多头投资者的"盈亏不对称"——一只股票从10元涨到100元，多头赚900\%，空头亏900\%（假设空头在10元开始做空）。如果股票从100元跌到10元，多头亏90\%，空头赚90\%。计算两种情景下的"不对称程度"并讨论。
\end{enumerate}
\end{exercises}

\xuehaiinline{
\textbf{GameStop的完整故事（2021年1月）：}\\
\textbf{背景：}GameStop（美国一家视频游戏零售商）基本面持续恶化（数字下载取代实体光盘），大量对冲基金做空该股票——空头比率一度超过140\%。\\
\textbf{触发：}Reddit社区WallStreetBets的用户发现这个极度做空的情况——开始集体买入看涨期权和现货。\\
\textbf{Gamma Squeeze的机制：}大量买入Call→做市商被迫买入股票来对冲Δ→股价上涨→Call变为实值→做市商需要买入更多股票对冲→进一步推高股价→恶性循环。\\
\textbf{结果：}股价从20美元涨到最高483美元。Melvin Capital（空头主力）损失约70亿美元被迫关闭。做市商Citadel Securities注资20亿美元救助Melvin。Robinhood等券商因保证金要求激增而暂时禁止用户买入GameStop（引发了巨大争议）。\\
\textbf{监管后果：}SEC在2021年10月发布了GameStop事件报告——结论是市场结构没有根本性问题，但建议全面审查"订单流付费"（PFOF）——零佣金券商通过将客户订单卖给做市商盈利的商业模式的合理性。2024年SEC已经提出了禁止PFOF的提案——但争议巨大。\\
\textbf{争议的核心问题：}"散户联合逼空机构"是资本市场民主化的胜利，还是市场操纵的新形式？做市商的"做市豁免"（允许做市商在市场波动时暂停对某些股票提供流动性）是否被滥用了？
}

\newpage
% ============================================================
\chapter{GameStop、Meme股票与散户化交易}
% ============================================================

内容同上一课的学海拾遗已经覆盖了GameStop的核心——但作为独立成课，本文从"制度冲击"的角度深入分析。

\begin{scenario}
\textbf{零佣金的革命——Robinhood如何改变了资本市场}

2015年，Robinhood推出了零佣金股票交易——颠覆了华尔街的经纪商收费模式。传统券商每笔交易收费5-10美元时，Robinhood说"免费"。3年后，Interactive Brokers等全线跟进零佣金。但Robinhood的创新不止于零佣金——它还让期权交易变得像按按钮一样简单。2021年，Robinhood超过2000万用户——其中约一半是第一次参与股票交易。Robinhood的商业模式是"订单流付费"（PFOF）——将客户订单卖给高频做市商（如Citadel Securities），从中赚取回扣。这是一个有争议的模式：客户是否得到了最好的执行价格？
\end{scenario}

\begin{qualitative}
\textbf{Meme股票现象的制度背景}

\begin{enumerate}
  \item \textbf{零佣金革命的副作用：}交易成本趋近于零鼓励了高频交易和投机行为——尤其是小金额、短持有期的"赌一把"型交易
  \item \textbf{社交媒体的"羊群效应"加速器：}WallStreetBets有超过1000万订阅者。Reddit帖子的传播速度远快于传统金融信息（分析师报告、新闻）——信息传播的基础设施变了
  \item \textbf{零保证金期权交易：}Robinhood允许甚至鼓励零售投资者交易期权——而期权是一种复杂的高杠杆工具。一个18岁、没有收入、没有投资经验的人可以在一分钟内下单买入带有无限损失的裸期权
  \item \textbf{"FOMO"（害怕错过）驱动的定价：}Meme股票的价格不再反映基本面——它反映的是社交媒体上的"情绪动量"
\end{enumerate}
\end{qualitative}

\begin{quantitative}
\textbf{AMC Entertainment的案例}

2021年初，AMC股价约2美元。6月峰值达到约72美元——涨了36倍。但AMC的财务数据在这期间恶化——2021年第一季度亏损5.7亿美元，收入同比下降84\%。基本面和技术面产生了极端背离。

\textbf{订单流付费（PFOF）的隐性成本：}SEC的2018年开始实施的"最佳执行"规则要求券商在可能的情况下为客户争取最好的价格。但当客户订单被卖给Citadel等做市商时——做市商在订单流上赚取了买卖价差（可能远高于公开市场上的价差）。学术界对PFOF的影响有不同的估计——一些人认为PFOF让零售投资者每年"损失"数十亿美元；另一些人认为做市商执行的"价格改善"弥补了PFOF的隐性成本。
\end{quantitative}

\begin{lessonsummary}
\begin{enumerate}
  \item 零佣金交易和社交媒体从根本上改变了资本市场的价格形成机制——散户的力量从未如此之大
  \item Meme股票现象暴露了两个根本性问题：市场效率假说在社交媒体时代还成立吗？散户期权交易的风险是否超出了个体可以承受的范围？
  \item PFOF（订单流付费）的商业模式存在根本性的利益冲突——客户的利益和平台的利益不完全一致
  \item 中国未来的监管者需要借鉴美国GameStop事件的教训——散户化交易在任何市场都是双刃剑
\end{enumerate}
\end{lessonsummary}

\begin{discuss}
\begin{enumerate}
  \item "金融民主化"的边界在哪里？Robinhood使低收入者可以参与市场——这是好事。但当这些参与者用借来的钱（或零保证金的期权）交易极其高风险的资产时——这是"民主化"还是"对金融风险的无知导致的赌博"？
  \item Reddit上的"联合买入"行为是否构成市场操纵？SEC的结论是——不构成，因为一群人各自做出独立的买入决策并公开讨论——这是"言论自由"保护的范围。但如果有一小群人私下协调行动呢？
  \item 中国A股市场的"散户化"特征——A股散户交易量占比约60-70\%（vs 美股约20\%）。如果Meme股票现象在中国发生，监管者的政策响应可能是什么？——中国市场对"市场操纵"的执法力度更严格、监管者可以直接暂停交易。
\end{enumerate}
\end{discuss}

\begin{exercises}
\begin{enumerate}
  \item 假设AMC在2021年1月的空头比率（Short Interest / Float）为80\%，日均交易量5000万股。如果股价从2美元涨到72美元——计算做空者的累计损失（假设平均做空价格5美元、做空股数1亿股）。这个损失在什么情况下会触发"轧空"的自我加强循环？
  \item 分析"Gamma Squeeze"的量化机制：一份看涨期权合约的行权价为20美元，股价从20美元涨到40美元，期权的Δ从0.5增加到0.9。如果做市商卖出了1万张这样的期权（每张覆盖100股）——当股价从20涨到40时，做市商需要买入多少股来对冲初始和最终的Δ？这对股价有什么影响？
  \item 你作为监管者，会如何处理PFOF？——完全禁止（SEC提案）、加强信息披露（要求券商披露平均价格改善）、或维持现状？为你的选择辩护。
\end{enumerate}
\end{exercises}

\xuehaiinline{
\textbf{2021年的"Meme股"时间线：}\\
1月11日：Ryan Cohen（宠物电商Chewy的联合创始人）加入GameStop董事会——WSB热情高涨。\\
1月13日：GameStop股价从20美元涨到31美元。\\
1月22日：股价涨到65美元。Melvin Capital宣布关闭GameStop空头仓位。\\
1月27日：特斯拉CEO Elon Musk发推"Gamestonk!!"——股价涨到347美元。\\
1月28日：Robinhood暂停用户买入GameStop和其他Meme股票（因存款要求激增至30亿美元）——引发了巨大争议。\\
1月29日：股价达到峰值483美元。\\
2月之后：股价回落至40-50美元区间——但仍在基本面价值（约10-20美元）之上。\\
后续：GameStop在2021年4月通过高位增发（At-the-Market Offering）融资了约17亿美元——成了GameStop事件最大的赢家。Robinhood的IPO招股说明书中将GameStop事件列为重大风险因素，并在2024年被SEC罚款了约4500万美元（因PFOF披露不完整）。
}

\newpage
% ============================================================
\chapter{2023年硅谷银行——数字时代的银行挤兑}
% ============================================================

\begin{scenario}
\textbf{Tweet触发的银行挤兑——420亿美元在几小时内消失}

2023年3月9日，硅谷银行（SVB）的储户通过网上银行系统在几个小时内取走了420亿美元——占该银行存款总额的约25\%。这是历史上最快的银行挤兑——不是排队的储户，不是打电话——而是一系列Tweet和WhatsApp消息引发了"数字挤兑"。仅仅24小时后，SVB被FDIC（美国联邦存款保险公司）接管。一个资产2118亿美元、被认为是"稳健"的银行，在48小时内倒塌。SVB不是2008年那些持有次贷产品的银行——它持有的是美国国债和机构MBS——理论上"最安全"的资产。
\end{scenario}

\begin{qualitative}
\textbf{SVB资产负债表错配的技术解剖}

\textbf{资产负债表的"死亡螺旋"：}
\begin{enumerate}
  \item 2020-2021年科技繁荣期，大量创业公司存款涌入SVB（存款从650亿增至1890亿）
  \item SVB用这些存款购入了大量长期（5-10年）美国国债和MBS——分类为"持有至到期"（HTM），按会计准则不需要按市值计价
  \item 2022年美联储开始加息500bp——长期债券价格暴跌（久期越长的债券跌得越惨）
  \item SVB的HTM组合的市值损失约180亿美元——但资产负债表上没有体现
  \item 2023年3月，SVB宣布需要出售210亿美元的AFS（可供出售）债券并融资22.5亿美元来弥补亏损
  \item 融资消息传出后，VC机构开始建议被投公司撤出SVB存款——Tweets和WhatsApp群组传播恐慌
  \item 存款人在几小时内取款420亿美元——SVB被迫紧急出售HTM资产——按市值变现的损失使其资本金不足以覆盖
  \item FDIC在3月10日接管
\end{enumerate}
\end{qualitative}

\begin{quantitative}
\textbf{SVB的久期缺口计算}

\begin{itemize}
  \item 总资产约2118亿美元，其中HTM证券约913亿（久期约6.2年）+ AFS证券约264亿（久期约3.7年）
  \item 加权平均资产久期：(913×6.2 + 264×3.7) / 1177 ≈ 5.6年
  \item 负债端——主要是活期存款（久期≈0）——但部分存款有固定期限
  \item 加权平均负债久期——约1.2年
  \item 久期缺口 ≈ 5.6 - (2118/2118) × 1.2 ≈ 4.4年
  \item 利率上升500bp时，资产市值变动 ≈ -5.6 × 5\% = -28\% ≈ -593亿美元
  \item 负债市值变动 ≈ -1.2 × 5\% = -6\% ≈ -127亿美元
  \item 净资产变动 ≈ -593 + 127 = -466亿美元
  \item SVB总股本约160亿美元——466亿的潜在亏损远超过全部股权
\end{itemize}

\textbf{关键点：}以上计算假设所有HTM资产都按市值卖给第三方——但实际上HTM可以持有至到期（如果不需要变现）。SVB的灾难性在于——它被挤兑了，被迫在不利条件下变现HTM——把账面损失变成了实际损失。
\end{quantitative}

\begin{lessonsummary}
\begin{enumerate}
  \item SVB倒闭的核心原因——极端期限错配（借短投长）+ 存款过于集中（单一行业）+ 社交媒体引爆了挤兑
  \item HTM（持有至到期）会计分类帮助SVB"隐藏"了180亿美元的亏损——直到挤兑逼迫它变现
  \item 数字时代的银行挤兑可以在几小时内发生——没有"冷静期"，没有"排队"——监管框架还没有适应这个新现实
  \item 美联储在SVB之后推出了银行期限资金计划（BTFP）——允许银行按面值向美联储质押债券借款——解决了"好资产、差时机"的流动性问题，但有道德风险
\end{enumerate}
\end{lessonsummary}

\begin{discuss}
\begin{enumerate}
  \item HTM会计分类是否应该取消？支持理由：它允许银行隐藏真实的利率风险；反对理由：如果取消，所有银行都必须按市值计价——债券价格波动将直接冲击银行利润表——在危机时可能加剧恐慌（"被强制Mark-to-Market加速死亡"）。
  \item 如果美联储有BTFP（按面值质押借款）作为流动性工具——SVB是不是本来可以避免倒闭？——是的，如果SVB在发现流动性问题之前通过BTFP借入了足够的资金来应对挤兑。但问题在于——SVB发现自己需要借钱的时候，挤兑已经发生了。流动性工具只有在"提前使用"时有效——在挤兑发生时，任何工具都只能减缓、而不能阻止恐慌的蔓延。
  \item 中国银行体系的存款保险制度——每个存款人最高偿付限额50万元人民币。如果发生类似SVB的事件，中国监管者的工具箱里有什么？——（1）央行可提供常备借贷便利（SLF）和中期借贷便利（MLF）向银行注入流动性；（2）存款保险基金管理公司可以提前偿付；（3）大型国有银行可以"托管"问题银行。但中国的工具是否能应对一个"数字挤兑"（速度远超传统挤兑）——没有经过实战检验。
\end{enumerate}
\end{discuss}

\begin{exercises}
\begin{enumerate}
  \item 假设一家中国中小银行持有大量长期债券（久期8年），存款主要由企业存款构成（久期≈0），银行资本充足率12\%。如果市场利率上升200bp——这家银行的久期缺口会使其净资产变动多少？是否可能突破资本金？
  \item 分析"未实现损失/资本金比率"——如果一个银行的未实现损失（Unrealized Loss）超过了其一级资本（Tier 1 Capital），虽然账面仍然"有偿付能力"——但实际上已经"经济上"破产了。2023年SVB的这个比率是120\%（未实现损失180亿/Tier 1资本150亿）。在2023年3月的美国银行业中，有多少家银行的这个比率超过了100\%？——根据一项研究，约有200家美国银行。
  \item 如果你是中国某地方银行的CFO，在2024年利率持续走低的背景下，你会如何管理利率风险？——延长还是缩短久期？为什么？
\end{enumerate}
\end{exercises}

\xuehaiinline{
\textbf{如果一切重来——SVB倒下的七个可干预节点：}\\
\textbf{节点1（2021年）：}SVB可以选择在其投资组合中加入更多的浮动利率资产（久期更短、利率风险更小）——但它没有。\\
\textbf{节点2（2022年）：}当美联储开始加息时，SVB可以选择对HTM资产进行"套期保值"（通过利率互换把固定利率转换为浮动利率）——但它没有。\\
\textbf{节点3（2022年末）：}SVB可以发行更多的长期债务来缩短资产负债表的久期缺口——但它没有。\\
\textbf{节点4（2023年初）：}SVB可以提前出售部分HTM资产并锁定损失、补充资本金——但它"等待情况好转"。\\
\textbf{节点5（2023年3月8日）：}当评级机构开始审查SVB时，管理层可以立即启动融资方案——而不是等到3月9日才宣布。\\
\textbf{节点6（2023年3月9日上午）：}在第一批Twitter消息发出后，SVB可以发布澄清声明、甚至提前关闭线上提款通道——但监管机构不允许这样做（因为不符合公平对待所有储户的要求）。\\
\textbf{节点7（2023年3月9日下午）：}当420亿美元的提款请求涌入时——SVB可以向美联储申请紧急贷款——但它的资产已经处于"恐慌性抛售"中。\\
\textbf{教训：}问题不是"为什么最后一个节点失败"，而是"为什么在前6个节点什么都没有做"。
}

\newpage
% ============================================================
\chapter{中国金融风险全景图}
% ============================================================

\begin{scenario}
\textbf{恒大——中国史上最大的企业违约案}

2021年12月，恒大集团——中国第二大房地产开发商——正式违约一笔2.6亿美元的离岸债，开启了其2.4万亿人民币总负债的重组之路。恒大的债务涉及数百家银行、信托公司和数十万购房者。它的倒下不是孤立的——碧桂园（2023年）、融创、世茂等相继违约。2024年中国房地产行业的债务违约累计超过1万亿人民币。这不是流动性危机——这是中国房地产"高周转+高杠杆"商业模式的系统性崩盘。
\end{scenario}

\begin{qualitative}
\textbf{中国金融风险的四大领域}

\textbf{领域一：房地产——"高杠杆"模式的终结}

开发商的平均杠杆率在2019年超过80\%（净负债/权益）。恒大2.4万亿总负债中，约40\%是银行贷款、30\%是应付账款（欠施工队和供应商的钱）、10\%是债券（含离岸债）、20\%是其他——"明股实债"的不透明债务结构难以估量。2020年8月央行推出"三条红线"（剔除预收款后的资产负债率>70\%，净负债率>100\%，现金短债比<1倍）——从政策出台到恒大违约只有14个月。

\textbf{领域二：地方政府隐性债务——"信号"被忽视}

城投公司通过发行"准政府债券"融资，投资者普遍认为它们有地方政府（最终是中央政府）的"刚性兑付"背书。总规模估计约40-60万亿——远高于地方政府公布的法定债务余额（约35万亿）。2023年贵州、云南、辽宁等地的城投公司出现逾期——刚性兑付的信仰在动摇。

\textbf{领域三：影子银行——从100万亿到60万亿的收缩}

2017年影子银行巅峰期接近100万亿。资管新规（2018年）强制打破刚兑、净值化管理、限制多层嵌套——影子银行规模到2024年压缩至约60万亿。但这个过程中，非标融资的真空没有被银行信贷和债券市场完全填补。

\textbf{领域四：中小银行风险——"碎片化"的脆弱性}

中国有约4000家银行——其中绝大多数是资产规模不足百亿的农商行和村镇银行。2022年河南和安徽的村镇银行"取款难"事件暴露了这些银行的风险——IT系统外包给第三方、大额关联贷款、异地存款的过度集中。
\end{qualitative}

\begin{quantitative}
\textbf{杠杆率的国际比较}

中国非金融部门总债务/GDP约280\%——与发达国家相当。但结构显著不同——美国企业债务占GDP的约80\%，中国约150\%。中国的杠杆率在过去十年增长极快——从2010年的约170\%到2024年的约280\%。

\textbf{"明股实债"的结构性风险：}开发商表面的"少数股东权益"占总权益的比例从2010年的约15\%上升到2023年的约35\%——这些"权益"中有相当一部分实际上是债（固定收益率、保本承诺、回购协议）。如果把这部分"隐藏的债"计入实际债务——开发商的真实净负债率可能比账面高出30-50个百分点。

\textbf{地方债的可持续性：}当一个地区的GDP增速大幅放缓（从10\%到3\%）、土地出让收入下降（从2千亿到5百亿）、利息支出增速超过收入增速——债务可持续性迅速恶化。简单的"债务/GDP"指标掩盖了利率和增速的交互作用——$(\text{债务}/\text{GDP})$的变化率 ≈ 利率 - 名义GDP增长率。当名义GDP增长率低于债务平均利率时（如2024年的中国部分地区），债务/GDP比率会自动上升。
\end{quantitative}

\begin{lessonsummary}
\begin{enumerate}
  \item 中国金融风险的四个主要领域——房地产、地方隐性债务、影子银行、中小银行——相互关联
  \item "明股实债"和"刚兑信仰"使真实风险被系统性低估
  \item 中国金融风险的特点是"高集中度"但具有制度性的"兜底"（中央政府信用）——这既降低了即时的系统性崩溃风险，也增加了长远的道德风险
  \item 解决中国金融风险的根本出路在于建立市场化的风险定价机制——让风险反映到价格中、让违规者承担后果
\end{enumerate}
\end{lessonsummary}

\begin{discuss}
\begin{enumerate}
  \item "中国会爆发系统性金融危机吗？"——主流的共识是"不"（因为中央政府有资产负债表空间兜底），但挑战性的观点是"这次可能不同"（因为地方财政、房地产、人口结构三重压力叠加）。用自己的分析框架给出判断，并指出你判断中最不确定的变量是什么。
  \item "道德风险"的困境——每一次救市（如2023年增发1万亿特别国债支持地方）都在强化"刚兑信仰"——下次危机来临时，政府将面临更大的压力。但如果不救——可能会触发连锁违约。这是"时间不一致性"问题的经典表现——最优的短期选择（救市）与最优的长期规则（不救）不兼容。
  \item 中国房地产周期和美国次贷周期的异同——相同点：低利率催生泡沫、杠杆积累、最终价格下跌。不同点：中国没有"次贷"（首付比例高）、中国开发商高杠杆(而非居民)、政策调控主动性更强。这些差异是"防火墙"还是"风险转移"（从居民端转移到开发商端）？
\end{enumerate}
\end{discuss}

\begin{exercises}
\begin{enumerate}
  \item 简化分析恒大资产负债表：总资产2.8万亿、总负债2.4万亿。拆解负债端——银行贷款约9600亿、债券约2400亿（含离岸债）、应付账款约7200亿、其他约4800亿。资产端——土地储备约1.2万亿（假设已减值30\%）、在建项目约1万亿（假设折旧20\%）、现金约500亿、其他约500亿。计算恒大的实际净资产——如果按"真实价值"而不是账面价值估算。
  \item 用"债务/GDP"的数学公式分析——如果一个省份的名义GDP增速从10\%降到3\%，地方债务平均利率从4\%降到3\%——债务/GDP比率是上升还是下降？请给出计算和结论。
  \item 从中国人民银行或国家金融监管总局官网下载最新的《中国金融稳定报告》，找到关于"中小银行风险"的章节。分析监管层认为最突出的风险点是什么，并提出你的不同意见（如果有）。
\end{enumerate}
\end{exercises}

\xuehaiinline{
\textbf{1994年分税制改革——地方债问题的制度根源：}1994年之前，地方政府的财政收入占全国财政收入的约70\%——与它承担的支出责任基本匹配。1994年分税制改革后，增值税的75\%上缴中央、企业所得税按隶属关系划分——地方政府的财政自主收入大幅下降（到2024年约45\%），但支出责任没有减少（甚至增加了）。地方政府通过"土地财政"（卖地收入归地方、不进入一般公共预算）弥补了缺口。当房地产市场下行导致土地出让收入暴跌（2021年全国土地收入约8.7万亿，2023年约4.7万亿）——地方政府的收入来源出现了巨大缺口。这个缺口被隐性债务（通过城投公司）填补。\\
\textbf{解决方案的几种路径：}（1）重新调整中央和地方的财权和事权分配——增加地方政府的自主财力；（2）建立地方税体系——房产税和消费税是两个候选——但政治阻力巨大；（3）硬化地方政府预算约束——允许地方政府破产（中国《预算法》不允许，但部分学者呼吁建立"地方财政重整"制度）——但政治风险同样巨大。
}

\newpage
% ============================================================
\chapter{数字货币与DeFi——范式转移还是金融闹剧？}
% ============================================================

\begin{scenario}
\textbf{FTX的80亿美元黑洞——"革命"如何变成了"闹剧"}

2022年11月，FTX——全球第二大加密货币交易所、估值320亿美元——在几天内崩溃。其创始人Sam Bankman-Fried（SBF）被美国司法部指控七项欺诈罪名，2023年12月被定罪，面临最高100年以上监禁。FTX的崩溃揭示了整个加密货币行业的问题：客户资金被挪用（Alameda Research——SBF的另一家公司——从FTX客户账户借走了80亿美元用于高风险交易）、没有独立董事会、没有外部审计、没有合规体系、创始人自己决定钱的用途。加密世界的核心叙事——"去中心化""透明""信任less"——在FTX这个"中心化"交易所身上完全不适用。
\end{scenario}

\begin{qualitative}
\textbf{区块链的核心思想}

区块链是一种分布式账本技术——数据不存储在单个中心服务器上，而是由网络中所有参与者共同维护。每一次交易被打包成一个"区块"，按时间顺序连接成"链"——因此无法被单方面篡改。

\textbf{主要应用场景：}
\begin{itemize}
  \item \textbf{加密货币（Bitcoin, ETH, etc.）：}去中心化的数字"货币"——不需要银行作为中介
  \item \textbf{智能合约（Smart Contract）：}在区块链上运行的程序——自动执行合约条款（如"A收到付款后自动向B转移数字资产所有权"）
  \item \textbf{DeFi（去中心化金融）：}基于智能合约的链上金融服务——借贷、交易、抵押、保险——没有中介机构
  \item \textbf{NFT（非同质化代币）：}代表数字资产所有权的唯一性代币——可以用来"证明"一段视频、一张图片的"原始作者"
\end{itemize}
\end{qualitative}

\begin{quantitative}
\textbf{比特币的供应量限制}

比特币的程序代码中写着——总供应量永远不会超过2100万个。目前（2024年）已经挖出了约1950万个（约93\%）。每四年"减半"一次——区块奖励从最初的50 BTC降至目前的3.125 BTC。到2140年左右将开采完最后一点比特币。

\textbf{PoW的能耗：}比特币网络每年消耗约100-150 TWh的电力——相当于一个小型发达国家的年用电量。每笔比特币交易的能耗约700 kWh——相当于一个普通美国家庭约三周的用电量。这是比特币在ESG方面受到的最大批评。

\textbf{数字人民币（e-CNY）——全球最大的央行数字货币：}与比特币完全相反——e-CNY是完全中心化的（央行控制供应、央行记录交易）。到2024年累计交易额超过2万亿人民币。优势在于：降低现金运营成本、增强反洗钱能力、增强货币政策传导效率（未来可能的"直升机撒钱"——直接给每个公民的数字钱包发钱）。
\end{quantitative}

\begin{lessonsummary}
\begin{enumerate}
  \item 加密货币是一项有历史意义的实验——但它至今没有证明自己是比传统金融更好的支付系统或价值储存工具
  \item "去中心化"的叙事和现实之间存在巨大差距——大多数加密货币交易在中心化交易所（CEX）上完成，这些交易所的风险和传统金融机构一样
  \item 数字人民币（e-CNY）代表了一条完全不同的道路——在现有金融体系的基础上用数字技术创新效率，而不是"推倒重建"
  \item 加密货币的未来取决于两个变量：监管政策的走向（美国SEC的执法、欧盟MiCA法规）和技术突破（可扩展性、用户友好性）
\end{enumerate}
\end{lessonsummary}

\begin{discuss}
\begin{enumerate}
  \item 2024年美国SEC批准了现货比特币ETF（贝莱德等发行），这是否标志着"比特币作为资产类别"被主流金融接纳？——是：ETF让传统投资者可以像买卖股票一样买卖比特币；不是：比特币作为"货币"的功能仍然极其有限（支付速度慢、成本高、波动大）。
  \item FTX毁灭后，"监管加密货币"在全球范围内成为共识。但"如何监管"仍有巨大分歧——是"像证券一样监管"（信息披露、注册、反欺诈），还是"像商品一样监管"（美国CFTC的立场），还是"银行的监管"（欧盟MiCA法规的立场）？
  \item "区块链可能改变的不只是金融"——去中心化身份（DID）、供应链溯源、版权保护、政务数据共享——这些应用场景可能比加密货币更有价值。你对区块链最看好的非金融应用是什么？
\end{enumerate}
\end{discuss}

\begin{exercises}
\begin{enumerate}
  \item 计算一个比特币挖矿的盈利经济——假设矿机（蚂蚁S19）功耗3250W、效率30J/TH、算力100TH/s、全网算力400EH/s、电费0.4元/度、区块奖励3.125 BTC、当前比特币价格40万人民币。计算单台矿机的日收入、日电费、日净利润。回本周期大约是多少？
  \item 分析e-CNY对现有支付体系（支付宝/微信支付）的影响——e-CNY和支付宝的五个核心区别是什么？e-CNY能否"取代"支付宝？——大概率不能，因为e-CNY定义的是"钱本身"（M0），而支付宝定义的是"钱的通道"——两者不是同一个层次。
  \item 选择一个DeFi协议（如Uniswap——去中心化交易所），分析它的运作机制——流动性池（Liquidity Pool）vs 传统订单薄（Order Book）——两者的优缺点。在什么情况下DeFi优于CeFi（中心化金融）？
\end{enumerate}
\end{exercises}

\xuehaiinline{
\textbf{FTX崩溃的完整时间线（2022年11月）：}\\
11月2日：CoinDesk（加密货币新闻网站）发布文章——披露Alameda Research持有的最大单一资产是FTX的交易所代币FTT（约50亿美元）——Alameda的资产负债表极度不透明。\\
11月6日：Binance CEO赵长鹏宣布将清算Binance持有的全部约5.8亿美元FTT（竞争对手的代币）。\\
11月7-8日：FTX客户开始大规模撤资——三天内撤资约60亿美元。FTX暂停提款。\\
11月9日：Binance宣布"意向收购"FTX（非约束性）——几个小时后放弃。\\
11月11日：FTX美国分公司宣布破产。约100万债权人和投资者受影响。\\
11月12日：SBF辞职。发现客户资金缺口约80亿美元。\\
12月12日：SBF在巴哈马被捕。\\
2023年11月：SBF被定罪，七项罪名全部成立。\\
2024年3月：SBF被判处25年监禁。\\
\textbf{关键问题：}FTX客户的所有资产——理论上应该是"1:1"持有的（每个客户存入1个比特币，FTX就应该持有1个比特币）——但实际上FTX将客户资产借给了Alameda做交易。这是一个存在了几百年的金融犯罪——"挪用客户资产"——但FTX用"加密货币的创新"包装了它。
}

\newpage
% ============================================================
\chapter{量化交易与HFT——算法在管理市场吗？}
% ============================================================

\begin{scenario}
\textbf{2010年"闪崩"——9\%的暴跌在几分钟内发生了什么？}

2010年5月6日，美国股市在14:32到14:47之间的15分钟内暴跌了约9\%（道指跌了998点）——然后同样迅速地恢复了。这是历史上最大的单日盘中跌幅之一。原因是什么？事后调查报告（CFTC-SEC联合报告）指出：一个交易员通过算法卖出了41亿美元的E-Mini S\&P 500期货合约（他没有使用任何价格或时间限制的算法，仅仅使用了"成交量参与"算法）。当他的算法系统性地卖出时，高频交易商的做市算法开始"竞赛性"地撤单——因为订单不平衡信号触发了它们的风险控制——做市商集体退出市场，造成了瞬间流动性枯竭——价格在没有买家的情况下"坠崖"。算法不是"恶意"的——只是"相互影响"产生了谁都没有预料到的结果。
\end{scenario}

\begin{qualitative}
\textbf{三类量化交易策略}

\begin{enumerate}
  \item \textbf{高频交易（HFT）：}利用极快的计算机和低延迟网络在微秒级别进行交易——做市商在买卖价差之间获利、套利者在不同交易所之间的价格差异中套利
  \item \textbf{统计套利（Stat Arb）：}基于历史统计关系的交易策略——当两类资产之间的价差暂时偏离历史范围时，做多被低估的、做空被高估的
  \item \textbf{动量/反转策略：}基于趋势跟踪——"追涨杀跌"的系统化版本
\end{enumerate}

\textbf{HFT的社会争议：}
\begin{itemize}
  \item 支持：HFT降低了买卖价差（学术研究表明约50-70\%）→降低了所有投资者的交易成本
  \item 反对：HFT在毫秒级别比人类投资者更快地看到订单流——这构成一种'不公平的隐性信息优势"
\end{itemize}
\end{qualitative}

\begin{quantitative}
\textbf{做市商的买卖报价策略}

最优报价宽度（Optimal Bid-Ask Spread）近似公式：
\[
\text{Spread} \approx \text{Inventory Cost} + \text{Adverse Selection Cost}
\]

做市商面临的"逆向选择"问题——当有人比你更了解股票的真实价值时，你（做市商）在和"知情交易者"交易时必然亏钱。因此做市商通过买卖价差来弥补这种损失。HFT通过更快的速度和更先进的风险管理可以降低Adverse Selection Cost——从而为所有市场参与者提供更窄的价差。

\textbf{中国量化交易的DMA策略（2024年）：}DMA（多空收益互换）是一种结构型产品——投资者支付一定费用，获得一个"多空组合"的总收益。2024年2月A股大跌时，大量DMA策略的底层空头头寸（用股指期货做空）因为现货下跌而"多"的部分巨亏——投资者无法追加保证金——券商被迫平仓——进一步加剧了市场下跌。一套以"市场中性"为名义的量化策略，在极端行情中变成了"市场极度非中性"的放大器。
\end{quantitative}

\begin{lessonsummary}
\begin{enumerate}
  \item 量化交易和高频交易已经占据了全球股市交易量的约60-70\%——不是"要不要"的问题，而是"如何管理"的问题
  \item HFT显著降低了交易成本，但增加了市场脆弱性——闪崩、速度竞赛、订单操纵
  \item 中国量化交易在2024年受到强监管——核心问题不是量化本身，而是监管框架的适应性
  \item 量化交易的最大社会价值不是"赚钱"——而是为市场提供了流动性和更有效的价格发现
\end{enumerate}
\end{lessonsummary}

\begin{discuss}
\begin{enumerate}
  \item "闪崩"是否意味着市场已经不安全了？——从另一个角度看，2010年闪崩在15分钟内自动恢复了——市场具有"自我修复"能力。2020年疫情冲击时，市场在几周内下跌了30\%——同样也恢复了。市场在极端情况下的"韧性"可能比我们想象的要强。
  \item "时间公平"（Time Fairness）的问题——HFT比人类快了几毫秒，这是否意味着市场存在根本性的不公平？如果两个人同时想买同一只股票，但HFT的交易算法比人类的指尖快10毫秒——它"应该"先成交吗？
  \item 中国在2024年引入的"程序化交易报备"制度——要求所有使用算法的交易者在监管机构注册。这能解决什么？不能解决什么？——能解决"监管不知道市场中有多少交易是算法执行的"；不能解决"算法之间的相互作用产生的系统性风险"。
\end{enumerate}
\end{discuss}

\begin{exercises}
\begin{enumerate}
  \item 模拟一个简单的做市策略：给定股票当前价格100元，假设该股票价格服从正态分布（标准差1元/天）。做市商报价：卖价100.10元、买价99.90元（价差0.20元）。每次成交100股，一天交易500次。假设每天股价随机变动（服从给定的分布），计算做市商一天的利润和库存风险。如果价差缩小到0.10元呢？
  \item 分析"闪电崩盘"的机制：假设市场上有一个大卖家在15分钟内卖出100亿份E-Mini合约（正常日成交量约200亿份）。流动性提供者的"订单簿"深度在价格下跌时的变化——如果每10个点的价格区间只有1亿的买单深度，剩下的99亿卖单需要价格下降多少才能被吸收？
  \item 用Python从Tushare下载A股1分钟高频数据，计算某日盘中"成交量的自相关"——即连续1分钟内的成交量是否有正相关？如果有，说明市场存在"羊群效应"的特征。
\end{enumerate}
\end{exercises}

\xuehaiinline{
\textbf{2010年Flash Crash的时间线（来自CFTC-SEC联合报告）：}\\
14:32：Waddell \& Reed（一家共同基金）的交易员开始执行一个"卖空41亿美元E-Mini S\&P 500"的套保程序——只用了"成交量参与"算法（VWAP），没有设置价格限制、没有设置时间限制。\\
14:32 - 14:41：该算法卖出了约19亿美元合约。HFT做市商开始累积空头头寸——它们通过做多其他股票来对冲——并在订单簿上撤单（减少流动性供应）。\\
14:41 - 14:44：做市商的净库存超过了其风险容忍度——开始系统性撤出市场（取消所有报单）。买单深度骤降至正常水平的约1\%。\\
14:45 - 14:47：流动性坍塌——价格在几乎没有买单的情况下，从约1060点跌到约986点（含一次"交叉"——没有交易、只有订单）。E-Mini价格一度跌停了约10\%。\\
14:47 - 15:00：价格在买单重新进入（包括CME的"止损断路器"——暂停交易5秒——让市场参与者重新评估）后迅速恢复。\\
\textbf{监管措施：}（1）引入"限价上下行机制"（Limit Up-Limit Down, LULD）——当股票价格在5分钟内变动超过一定比例时，暂停交易5分钟；（2）加强对"无价格限制算法"的交易监管；（3）要求所有算法交易者在CFTC注册。\\
\textbf{讽刺的结局：}Waddell \& Reed的交易员坚持认为他的交易"一直保持在正常市场条件下的合规交易"——而CFTC也无法证明他违反了任何明确的法律。这说明了——闪崩不是"一个人做错事"的结果，而是"系统在一组特定条件下的自动行为"。
}

\newpage
% ============================================================
\chapter{金融伦理——资本市场到底在服务谁？}
% ============================================================

\begin{scenario}
\textbf{瑞幸咖啡——22亿的财务谎言如何被揭穿}

2020年1月31日，做空机构浑水（Muddy Waters）发布了一份89页的报告——指控瑞幸咖啡（Luckin Coffee）伪造了约22亿人民币的交易额。浑水派出了92个调查员在1513家瑞幸门店观察和记录了981个"店日"的客流量——发现瑞幸公布的单店日均销量被夸大了约69\%。报告发布后瑞幸股价大跌，然后在两个月内继续运作——直到2020年4月2日瑞幸"自曝"（特别委员会确认了22亿造假），股价暴跌85\%。2020年6月，瑞幸从纳斯达克退市。2022年SEC对其罚款1.8亿美元。瑞幸成为中概股历史上最大规模的财务造假案之一。
\end{scenario}

\begin{qualitative}
\textbf{金融伦理的核心问题}

\begin{enumerate}
  \item \textbf{投资者保护 vs 资本形成：}资本市场的存在是为了让企业融资（降低融资成本），同时也是为了让投资者分享经济增长（获得合理回报）。当两者冲突时——比如企业造假融资——投资者保护应该优先
  \item \textbf{利益冲突的来源：}
  \begin{itemize}
    \item 委托-代理问题——经理人追求自身利益（奖金、期权、声誉）可能偏离股东利益
    \item 利益输送——大股东通过关联交易转移上市公司利益
    \item 信息不对称——内部人知道的远比外部人多
    \item 市场操纵——大资金通过虚假交易、散布假消息影响价格
  \end{itemize}
  \item \textbf{道德风险：}当一个人或机构（如银行）知道"出了问题有人兜底"的时候——会冒比正常情况下更大的风险。这是所有金融监管面临的终极难题
\end{enumerate}
\end{qualitative}

\begin{quantitative}
\textbf{内幕交易的实证研究}

内幕交易人在其交易后的6个月内获得了约3-5\%的超额收益（超越市场平均）。虽然3-5\%看似小，但乘以巨大的交易量和利用杠杆可变的非常可观。

\textbf{做空机构的社会价值：}浑水在2010-2024年间发布了约100份做空报告——其中最成功的是瑞幸（被退市）、辉山乳业（停牌退市、创始人被逮捕）、嘉楠科技（股价腰斩）。浑水的创始人Carson Block说："我们做的是财务审计——不是由被审计公司付钱、而是由我们自己付钱的审计。"这个"付费模式"的转换——从公司付费给审计师（传统模式，利益冲突明显）到市场付费给发现者（做空模式，激励机制正确）——是金融市场监督的一个有趣探索。

\textbf{ESG评级与投资回报：}MSCI的学术研究表明——高ESG评分公司的股本成本（WACC的一部分）平均比低ESG公司低约0.5-1.0\%。这不是说"做好事有超额收益"——而是说"做坏事的风险被市场定价了"。低ESG评级的公司面临更高的监管风险、诉讼风险和声誉损失风险。
\end{quantitative}

\begin{lessonsummary}
\begin{enumerate}
  \item 金融市场的核心矛盾——企业融资需求和投资者保护之间的平衡——永远没有完美答案
  \item 利益冲突是金融伦理的核心问题——每一个"创新"都可能创造新的利益冲突
  \item 做空机构在市场中扮演了"监管延伸"的角色——但本身同样存在利益冲突（做空后发布负面报告——先做空再发布报告，做空者从中获利）
  \item 绿色金融和ESG投资代表了金融"向善"的方向——但"漂绿"（Greenwashing）和ESG评级方法论的不透明使这个方向的实践仍然存在深层争议
\end{enumerate}
\end{lessonsummary}

\begin{discuss}
\begin{enumerate}
  \item 浑水做空瑞幸的模式——先做空、再发布报告——是否本身存在利益冲突？如果瑞幸的指控是错误的，浑水已经从中牟利了。但如果没有做空机构——瑞幸的造假可能要到很久之后才会被发现。这是"市场有效"的必要机制，还是"滥用做空的机制"？
  \item 2015年A股股灾中，一些上市公司大股东通过"高送转"（10送20）的消息在高位减持套现。技术上完全合规——但从伦理上怎么看？"合规但不合理"的问题如何解决？
  \item ESG评级机构的"评分标准"不透明——一家石油公司可能因为"治理好"获得ESG高分（如壳牌），而一家清洁能源创业公司可能因为"治理不成熟"获得低分。ESG评分驱动了数万亿美元的资本配置——但如果评分标准本身就是有问题的——资本的配置是否在被误导？
\end{enumerate}
\end{discuss}

\begin{exercises}
\begin{enumerate}
  \item 分析瑞幸咖啡的案例：瑞幸的商业模式（"烧钱换增长"——通过补贴快速扩张，目标是"熬到盈利的那一天"）本来就有极高的风险。财务造假的动机是什么？用"委托-代理问题"的框架分析——管理层（代理方）追求什么？股东（委托方）追求什么？两者在什么条件下产生了不可调和的冲突？
  \item 你是一个投资银行的分析师。你的投行为一家IPO公司做承销。你所在的分析部需发布对这家公司的"独立研究报告"。你会给出什么评级？如果给出"买入"有助于投行获得承销业务——你的决策会受到影响吗？这是否需要监管干预？
  \item 查询你感兴趣的一只股票（或ETF）的ESG评级（可通过MSCI的网站免费查询部分数据）。分析它的ESG评分和实际投资风险之间的关系。ESG评分是否反映了你对该公司的风险认知？
\end{enumerate}
\end{exercises}

\xuehaiinline{
\textbf{金融伦理的一个思想实验：}假设你是一个投资经理。你的一个客户是退休基金——管理着数万名退休教师的养老金。你的一个交易对手是一个"金融危机"中做空了你的客户所持有的股票的对冲基金经理。你会和这个对冲基金经理交易吗？如果交易能为你客户的退休基金盈利——你会做吗？\\
\textbf{金融学的答案：}会的——因为你的受托责任是为客户追求最佳风险调整后回报——和谁交易、谁在做空——不影响交易的价值。\\
\textbf{伦理的追问：}你知道这个对冲基金经理的盈利部分来自做空了你客户可能购买的同一只股票——而做空行为本身可能让这些公司的融资更加困难——最终影响了你客户（退休教师）的就业和退休金安全——你会改变你的交易吗？\\
\textbf{这个思想实验说明：}金融伦理不是一个"对或错"的问题——它是一个连续谱——"谁的利益优先"是每一个金融从业者每天都要面对的问题。CFA Institute（特许金融分析师协会）的《道德准则和职业行为标准》提供了这个领域的专业标准——将客户利益置于个人利益之上、将诚信置于短期利润之上。
}

\end{document}
